% SOPNet: A Signed-Onset Field Network for Automatic P-Wave First-Motion
% Polarity Determination.
%
% DRAFT. Target venue: Applied Geophysics (Springer).
% Numbers tagged \TBD{} are filled from outputs/paper/tables/*.csv once
% gates G1--G3 complete.
% NOTE: swap the document class to the journal template at submission time.

\documentclass[11pt]{article}
\usepackage[margin=2.5cm]{geometry}
\usepackage{amsmath,amssymb}
\usepackage{booktabs}
\usepackage{graphicx}
\usepackage{siunitx}
\usepackage[colorlinks=true,linkcolor=blue,citecolor=blue,urlcolor=blue]{hyperref}
\usepackage{xcolor}
\usepackage{natbib}

\newcommand{\TBD}[1]{\textcolor{red}{[#1]}}

\title{SOPNet: A Signed-Onset Field Network for Automatic P-Wave \\
First-Motion Polarity Determination}
\author{\TBD{author list}}
\date{\TBD{date}}

\begin{document}
\maketitle

\begin{abstract}
First-motion polarity of the P wave is the primary observable for constraining
earthquake focal mechanisms, and its automatic determination has become a
standard component of deep-learning seismology pipelines. Existing models,
however, are almost universally formulated as three-class softmax classifiers
(up / down / unknown), evaluated with heterogeneous datasets, splits and label
policies, which makes published numbers difficult to compare and provides
little insight into behaviours that matter in practice, such as noise and
arrival-time errors. We present SOPNet, a compact one-dimensional
encoder--decoder that instead predicts a \emph{signed onset field}: a
continuous function over the waveform window whose absolute maximum locates the
P onset and whose sign at that maximum gives the polarity. The field is
supervised only on records with known up/down labels using a weighted smooth
$L_1$ loss, and is regularised by a differentiable polarity-consistency loss
and a physical inversion-consistency constraint $f(-x) = -f(x)$. We assemble a
unified benchmark from five open datasets (SCSN, DiTing, INSTANCE, PNW, TXED;
7.78 million waveforms from 1.12 million events) with a shared,
event-level split, and we retrain five public deep-learning polarity models
(Ross, DiTingMotion, RPNet, EQPolarity, CFM) under their official input
protocols. On a strict up/down decision metric, SOPNet reaches
97.8\% accuracy, 0.978 macro-$F_1$ and 0.956 MCC on
327{,}231 known-polarity test waveforms, outperforming all five baselines,
with paired event-clustered bootstrap intervals excluding zero. Robustness
experiments on a fixed subset show graceful degradation under additive noise
(0.977 clean to 0.780 at \SI{-5}{dB}, the strongest model at every noise
level), while all six models remain critically sensitive to P-reference
shifts of a few tens of milliseconds; we discuss this accuracy--tolerance
trade-off as an inherent consequence of fixed-window polarity
classification. The signed-field
formulation yields a simple, well-calibrated and easily reproducible polarity
model, and the benchmark protocol is released for reuse.
\end{abstract}

\section{Introduction}
\label{sec:intro}

The polarity of the first P-wave motion recorded at a station --- whether the
ground first moves up or down --- is one of the oldest and most informative
observables in observational seismology. Together with the take-off angle it
constitutes the classical raw material of focal-mechanism determination
\citep{hardebeck2002hash}, and dense first-motion polarity datasets now feed
routine earthquake monitoring, stress-field studies and physics-based
earthquake forecasting. Historically, polarities were read manually from
seismograms, which is accurate but prohibitively labour-intensive at the scale
of modern, machine-learning-driven catalogs containing millions of events.
Automating this reading task has therefore received sustained attention.

Early automatic approaches relied on hand-crafted features of the onset
waveform, most commonly the sign of the first extremum or of a short-term
average. These rules degrade rapidly for emergent onsets and noisy traces, and
their error rates were difficult to quantify. With the advent of deep learning
in seismology, polarity determination was reformulated as a supervised
classification problem: \citet{ross2018} trained a convolutional network to
pick P arrivals and assign polarity labels on Southern California data, and a
series of recent models extended this line of work with deeper or more
specialised architectures, including DiTingMotion \citep{zhao2023ditingmotion},
RPNet \citep{han2025rpnet}, EQPolarity \citep{chen2024eqpolarity}, CFM
\citep{messuti2023cfm} and PolarCAP \citep{chakraborty2022polarcap}. Depending
on the model, the problem is posed as binary (up/down) or three-class
(up/down/unknown) classification, and reported accuracies are often above
\SI{95}{\percent} on held-out records.

Three gaps motivate this paper. First, the classifiers treat polarity as a
discrete label emitted by a softmax head, discarding the continuous, signed
nature of the first motion, and they provide no explicit, differentiable link
between the localisation of the onset and the polarity decision. Second,
comparability across publications is limited: models are trained on different
source datasets, with different preprocessing, different splits and different
label policies for records with ambiguous or unlabelled polarity, so that
cross-paper accuracy differences conflate model quality with protocol
differences. Third, robustness --- arguably the decisive property for field use
--- is rarely quantified: performance under ambient noise and, in particular,
under errors in the reference P arrival used to extract the analysis window is
seldom reported, even though a polarity model is only as trustworthy as its
behaviour when its input degrades.

We address these gaps with SOPNet, a polarity model built around a
\emph{signed onset field}, together with a strict benchmarking protocol. The
network is a compact one-dimensional convolutional encoder--decoder that maps
a waveform window to a continuous field $f(t)\in[-1,1]$. The absolute maximum
of $f$ marks the P onset; its sign at that maximum is the polarity. Because the
field target is constructed from the P arrival and the polarity label, the
model jointly learns localisation and polarity from a single regression
objective on records with known polarity. Two regularisers exploit the physics
of the representation: a polarity-consistency loss converts the field into a
differentiable polarity score, and an inversion-consistency constraint enforces
that reversing the input waveform reverses the field, $f(-x) = -f(x)$, a
symmetry that holds exactly for the first motion of an impulsive onset. At
inference, the magnitude of the field peak provides a natural confidence score
for optional rejection of ambiguous records and for probability calibration.

For evaluation, we assemble an open benchmark: five datasets (SCSN, DiTing,
INSTANCE, PNW and TXED) are harmonised into a single manifest with a shared,
event-level split that prevents leakage between correlated records of the same
earthquake. Records without a trustworthy polarity label act as unlabelled
data: they are excluded from the field-loss scope (``known-only'' training),
mirroring how such records are used in operational catalogs. Five public
baseline models are retrained from scratch --- no pretrained weights, to avoid
any form of leakage --- on the identical training split, each with its official
input protocol, and all models are compared with a strict up/down decision
metric that uses each model's own up/down score. Unknown-class predictions are
reported only as a supplementary rejection rate. We further evaluate all six
models on additive-noise sequences and on a grid of P-reference shifts, with
paired, event-clustered bootstrap intervals for the main comparisons.

The contributions of this work are:
\begin{itemize}
  \item a signed-field formulation of first-motion polarity, in which onset
  localisation and polarity are two readings of one continuous representation,
  trained with known-only supervision and regularised by polarity- and
  inversion-consistency losses;
  \item an open, event-level unified benchmark across five datasets and five
  retrained public baselines under official input protocols, with a strict
  up/down decision metric and per-source reporting;
  \item a quantitative robustness study of all six models under additive noise
  and P-reference shifts, with event-clustered paired bootstrap statistics;
  \item an ablation isolating the contribution of each component (binary
  classification, signed field, polarity-consistency and inversion-consistency
  losses) under identical training conditions.
\end{itemize}

The remainder of the paper is organised as follows.
Section~\ref{sec:related} reviews related work.
Section~\ref{sec:method} describes the signed-field model and its losses.
Section~\ref{sec:data} details the unified dataset, baselines and evaluation
protocol. Section~\ref{sec:results} reports main, per-source, ablation and
robustness results, and Section~\ref{sec:discussion} discusses their
implications and limitations.

\section{Related Work}
\label{sec:related}

\paragraph{First-motion polarity determination.}
\citet{ross2018} introduced the first widely adopted deep-learning polarity
model, a multi-branch convolutional network that picks the P arrival and
classifies its polarity on Southern California recordings; the same study
released the polarity labels that underpin several later datasets.
\citet{zhao2023ditingmotion} adapted a residual architecture with differential
input channels (DiTingMotion) and applied it to focal-mechanism inversion
\citep{zhao2023diting}. \citet{han2025rpnet} proposed RPNet, a residual network
designed for robustness across station conditions, and \citet{chen2024eqpolarity}
combined convolutional tokenisation with transformer blocks (EQPolarity) for
polarity determination and focal-mechanism applications. \citet{messuti2023cfm}
targeted volcanic and tectonic settings with a lightweight convolutional
feature model (CFM), and \citet{chakraborty2022polarcap} proposed an
encoder--decoder with a classification head (PolarCAP). These models share a
common formulation: the network output is a categorical distribution over
polarity classes, and the decision is an argmax. Our work keeps the
convolutional backbone family but replaces the categorical output with a
continuous signed field and derives the polarity from the field's own geometry.

\paragraph{Signed and continuous representations in seismology.}
Deep pickers such as PhaseNet \citep{zhu2019phasenet} and EQTransformer
\citep{mousavi2020eqtransformer} regress continuous phase-probability functions
over time and read arrivals from their maxima; this representation is now
standard for picking. For polarity, however, the corresponding idea --- a
continuous signed function whose sign encodes the first motion --- has not, to
our knowledge, been systematically evaluated against categorical classifiers,
nor combined with symmetry regularisation. Our signed onset field borrows the
localisation-by-maximum decoding of phase pickers and extends it with a sign
that is trained to match the first motion, so that a single output supports
both arrival refinement and polarity determination given a reference P time.
The inversion-consistency loss is inspired by the equivariance principle in
group-equivariant networks \citep{cohen2016group}: polarity inversion is an
exact symmetry of the underlying physical measurement, and we penalise its
violation directly.

\paragraph{Evaluation practice.}
Published polarity models report accuracy on their own splits; where unknown
labels exist, the treatment varies between ignoring them, folding them into a
third class, or excluding them at test time, and per-source behaviour is
rarely reported. Robustness analyses are similarly sparse, and we are not
aware of a six-model comparison under identical noise and arrival-error
protocols. The benchmark described in Section~\ref{sec:data} is intended to
fill this gap.

\section{Method}
\label{sec:method}

\subsection{Problem setting}
\label{sec:problem}

We consider vertical-component records for which a reference P arrival time
$t_\mathrm{ref}$ is available from a catalog or a picker, and we define the
task as determining the first-motion polarity given that reference --- the same
task posed to the baselines. Input waveforms are single-component windows of
$L$ samples centred on $t_\mathrm{ref}$ at a common sampling rate ($\SI{100}{Hz}$
in our experiments), band-pass filtered and amplitude-normalised
(Section~\ref{sec:data}). A polarity label $y\in\{+1,-1\}$ denotes an upward or
downward first motion; records with unknown or ambiguous polarity carry
$y=0$ and are used only as unlabelled context, never as field targets.

\subsection{Signed onset field}
\label{sec:field}

SOPNet is a one-dimensional convolutional encoder--decoder
(Figure~\ref{fig:architecture}). A multi-scale stem
with parallel convolutions feeds three downsampling stages (64, 96 and 128
channels), each followed by a residual block; the resulting feature map is
projected to 192 channels and processed by four dilated residual blocks
(dilation rates 1, 2, 4 and 8), giving a receptive field that spans the full
window. A symmetric decoder with skip connections restores the temporal
resolution, and a $1\times1$ convolution followed by a $\tanh$ produces the
signed onset field
\begin{equation}
  f_\theta: \mathbb{R}^{L} \rightarrow [-1,1]^{L}, \qquad
  f_\theta(x) = \big(f_\theta(x)_1,\dots,f_\theta(x)_L\big).
\end{equation}
The field is trained to be near zero away from the onset and to form a signed
lobe at the P arrival: the target $g$ is a Gaussian bump of width $\sigma$
(here \SI{10}{ms}) centred on the labelled P sample, scaled by the polarity
label, $g_i = y \exp(-(i-i_P)^2/(2\sigma^2))$. The model has
1.44 million parameters and operates on 400-sample windows,
matching the Ross/RPNet input convention.

\begin{figure}[t]
  \centering
  \includegraphics[width=\linewidth]{../../outputs/paper/figures/fig1_architecture.pdf}
  \caption{The SOPNet architecture. A multi-scale stem (parallel kernels of
  size 5/11/21) and a three-stage strided encoder with residual blocks feed a
  dilated bottleneck (rates 1, 2, 4, 8); a symmetric decoder with skip
  connections produces the full-resolution signed field
  $f(t)\in[-1,1]$. The P onset is read as the field maximum and the polarity
  as its sign; no classification head is used.}
  \label{fig:architecture}
\end{figure}

The field is decoded without any classification head. Let
\begin{equation}
  i^\star = \arg\max_i |f_\theta(x)_i|, \qquad
  s = f_\theta(x)_{i^\star}, \qquad
  c = |s| .
  \label{eq:decode}
\end{equation}
The polarity estimate is $\hat y = \mathrm{sign}(s)$ and $c\in[0,1]$ serves as a
confidence score. The peak position $i^\star$ refines the arrival time, but
under the fixed-centre training regime of this paper it should be read as an
internal consistency check rather than an independent picking capability
(see Section~\ref{sec:discussion}).

\subsection{Loss functions}
\label{sec:losses}

\paragraph{Weighted signed-field loss.}
Because the field target is zero almost everywhere, a plain mean-squared error
is dominated by the silent parts of the window. We weight the elementwise
smooth-$L_1$ loss by the target magnitude,
\begin{equation}
  \mathcal{L}_\mathrm{field}
  = \frac{1}{|\mathcal{K}|}\sum_{n\in\mathcal{K}}
    \frac{1}{L}\sum_{i=1}^{L}
      \big(1 + \beta\,|g^{(n)}_i|\big)\,
      \ell_{\delta}\!\big(f_\theta(x^{(n)})_i, g^{(n)}_i\big),
  \label{eq:fieldloss}
\end{equation}
where $\ell_\delta$ is the smooth-$L_1$ (Huber) function, $\beta=8$ up-weights
the onset neighbourhood, and $\mathcal{K}$ is the set of records with known
polarity. In the final recipe unlabelled records are excluded from
$\mathcal{L}_\mathrm{field}$ entirely (``known-only'' supervision); an
optional per-sample weight allows unlabelled records to be admitted with a
reduced coefficient, which we ablate in the supplement.

\paragraph{Polarity-consistency loss.}
To couple the field directly to the decision in Eq.~\eqref{eq:decode}, we form
a differentiable attention-weighted polarity score from the field itself,
\begin{equation}
  a_i = \frac{\exp(\kappa\,|f_\theta(x)_i|)}{\sum_j \exp(\kappa\,|f_\theta(x)_j|)},
  \qquad
  \hat s = \sum_i a_i\,f_\theta(x)_i,
  \label{eq:attn}
\end{equation}
and penalise sign disagreement with a softplus hinge,
\begin{equation}
  \mathcal{L}_\mathrm{pol}
  = \frac{1}{|\mathcal{K}|}\sum_{n\in\mathcal{K}}
    \mathrm{softplus}\!\big(-\gamma\, y^{(n)}\, \hat s^{(n)}\big),
  \qquad \gamma=1,\ \kappa=4 .
  \label{eq:polloss}
\end{equation}
As $\kappa\to\infty$, $\hat s$ approaches $f_\theta(x)$ at the field peak, so
minimising Eq.~\eqref{eq:polloss} sharpens the sign of the decoding quantity
itself rather than an auxiliary classifier.

\paragraph{Inversion-consistency loss.}
First-motion polarity is exactly equivariant under amplitude inversion: the
first motion of $-x$ is the opposite of that of $x$. We therefore penalise
deviations from $f_\theta(-x) = -f_\theta(x)$,
\begin{equation}
  \mathcal{L}_\mathrm{inv}
  = \big\| f_\theta(-x) + f_\theta(x) \big\|_1 ,
  \label{eq:invloss}
\end{equation}
applied on a random fraction of training batches. The total objective is
\begin{equation}
  \mathcal{L}
  = \mathcal{L}_\mathrm{field}
  + \lambda_\mathrm{pol}\,\mathcal{L}_\mathrm{pol}
  + \lambda_\mathrm{inv}\,\mathcal{L}_\mathrm{inv},
  \label{eq:total}
\end{equation}
with $\lambda_\mathrm{pol}=0.5$ and $\lambda_\mathrm{inv}=0.1$ in the final
recipe.

\subsection{Rejection and calibration}
\label{sec:calibration}

In operational use, records whose polarity cannot be determined reliably
should be rejected rather than mislabelled. The field peak magnitude $c$ in
Eq.~\eqref{eq:decode} provides a rejection score: a threshold $\tau$ is tuned
on the validation split to trade coverage for accuracy, and reported as a
supplementary selective-prediction curve. The same confidence is calibrated on
the validation split with Platt scaling \citep{platt1999calibration} --- a
one-parameter logistic map fitted to validation (confidence, correctness)
pairs --- reported via the expected calibration error (ECE) over
known-polarity records.

\section{Data and Benchmark Protocol}
\label{sec:data}

\subsection{Unified dataset}
\label{sec:dataset}

We harmonise five open datasets into one manifest:
SCSN (Southern California, with the refined polarity-rich catalog of
\citealp{ross2018,cheng2023focal}), DiTing \citep{zhao2023diting},
INSTANCE \citep{michelini2021instance}, PNW \citep{ni2023pnw} and TXED
\citep{chen2024txed}. All waveforms are resampled to \SI{100}{Hz}, band-pass
filtered between 1 and \SI{45}{Hz} (fourth-order Butterworth), and cut into
600-sample cache windows centred on the P pick; models consume the window
segment prescribed by their official input protocol
(Section~\ref{sec:baselines}). Polarity labels are mapped to the canonical
triad $\{+1, -1, 0\}$ (up / down / unknown). Table~\ref{tab:dataset}
summarises the corpus: 7.78 million waveforms from 1.12 million
events, split 80/10/10 at the \emph{event} level (events sharing an identical
waveform hash are kept together, preventing leakage through duplicated
records); the test split contains 777{,}720 waveforms and
111{,}588 events, of which 327{,}231 waveforms carry a known polarity
label. Unknown-polarity records are retained in all splits as unlabelled
context; they are never scored in the main metric.

\begin{table}[t]
  \centering
  \caption{Composition of the unified benchmark after harmonisation.
  ``Known U/D'' counts test waveforms with usable up/down labels; unknown
  records are excluded from scoring. Test events are counted over all test
  waveforms.}
  \label{tab:dataset}
  \small
  \begin{tabular}{lrrrrrr}
    \toprule
    Source & Waveforms & Train & Val & Test & Known U/D (test) & Test events \\
    \midrule
    SCSN     & 4{,}847{,}248 & 3{,}880{,}196 & 482{,}521 & 484{,}531 & 253{,}194 & 27{,}908 \\
    DiTing   & 2{,}734{,}748 & 2{,}186{,}864 & 273{,}031 & 274{,}853 & 64{,}048 & 78{,}835 \\
    PNW      &   183{,}909 &   151{,}336 &  15{,}906 &  16{,}667 & 9{,}742 & 4{,}392 \\
    TXED     &    10{,}692 &     8{,}533 &   1{,}020 &   1{,}139 & 143 & 181 \\
    INSTANCE &     4{,}964 &     3{,}988 &     446 &     530 & 104 & 272 \\
    \midrule
    Total    & 7{,}781{,}561 & 6{,}230{,}917 & 772{,}924 & 777{,}720 & 327{,}231 & 111{,}588 \\
    \bottomrule
  \end{tabular}
\end{table}

\subsection{Baselines}
\label{sec:baselines}

We include the five public baseline models available in the reference
implementation: Ross \citep{ross2018}, DiTingMotion \citep{zhao2023ditingmotion},
RPNet \citep{han2025rpnet}, EQPolarity \citep{chen2024eqpolarity} and CFM
\citep{messuti2023cfm}. Fair-comparison rules are as follows. (i) No pretrained
weights are used, for any model, including SOPNet; every network is trained
from scratch on the identical training split with the same optimiser schedule
(AdamW, cosine decay with warmup) and the identical augmentation policy
(additive noise at random SNR, polarity inversion, amplitude scaling). (ii) Each baseline consumes its official input: Ross and
RPNet 400-sample windows, DiTingMotion two channels with the official
derivative-sign second channel, EQPolarity 600 samples, and CFM its 160-sample
crop. (iii) The training objective and early-stopping monitor of each baseline
follow its official convention, with the monitor computed on up/down decisions
only so that the comparison is not distorted by the unknown class. (iv) All
models are evaluated on the same test samples, in the same order, with the
same metric.

\subsection{Metrics and statistics}
\label{sec:metrics}

The primary metric is the \emph{strict up/down decision}: for every test
waveform with a known polarity label, the model must commit to up or down
using its own up/down score (for categorical baselines, the binary
probability obtained by summing or restricting the up/down classes), and we
report accuracy, macro-averaged $F_1$ over the two classes, precision, recall
and MCC. For three-class baselines, predictions of ``unknown'' are treated as
abstentions and scored as errors in the strict metric; their rate is reported
separately as a rejection rate. SOPNet's threshold (Section~\ref{sec:calibration})
is likewise disabled in the strict metric and reported only as a selective
operating point. Because several waveforms may share an event and thus
correlated noise, pairwise differences against SOPNet are tested with a
paired bootstrap that resamples \emph{events} (not waveforms) with replacement
(1000 resamples), reporting the \SI{95}{\percent} percentile interval
of the macro-$F_1$ difference. Robustness experiments use a fixed
30{,}000-waveform subset of the test split (fixed seed) so that every
model sees exactly the same records; noise is injected at fixed SNR levels
(clean, 20, 10, 5, 0 and \SI{-5}{dB}) and the reference window is shifted on a
grid of $\{-0.20,-0.10,-0.05,0,+0.05,+0.10,+0.20\}$~s.

\section{Results}
\label{sec:results}

Figure~\ref{fig:examples} shows representative decisions of the full model:
two confident correct cases (clean and noisy onsets), one hard-but-correct
case at low confidence, and one high-confidence error, whose waveform
illustrates the emergent, low-amplitude onsets that dominate the residual
error.

\begin{figure}[t]
  \centering
  \includegraphics[width=\linewidth]{../../outputs/paper/figures/fig2_examples.pdf}
  \caption{Qualitative examples from the test set: waveform (top) and predicted
  signed field (bottom). Columns show two confident correct decisions, one
  low-confidence correct decision, and one high-confidence error. The red
  dashed line marks the reference P; the field peak (dot) locates the onset and
  its sign gives the polarity.}
  \label{fig:examples}
\end{figure}

\subsection{Main comparison}
\label{sec:main}

Table~\ref{tab:main} reports the strict up/down metrics on the known-polarity
test set (327{,}231 waveforms). SOPNet achieves 97.8\% accuracy,
0.978 macro-$F_1$ and 0.956 MCC, outperforming the five retrained baselines:
the best of them (CFM, 0.9727 macro-$F_1$) remains 0.005 behind, and the
largest gap (DiTingMotion) is 0.015. Paired event-clustered bootstrap
intervals over 63{,}044 events for the macro-$F_1$ difference against each
baseline are integrated into Table~\ref{tab:main}; all intervals exclude
zero. Rejection rates at each model's official decision rule are reported for
completeness in the supplementary table. SOPNet's validation-tuned rejection
threshold retains 99.5\% of known records and slightly raises their accuracy
(0.9806 versus 0.9782), while the calibrated confidence achieves an ECE of
0.031 (Figure~\ref{fig:calibration}).

\begin{table}[t]
  \centering
  \caption{Strict up/down results on the known-polarity test set
  (327{,}231 waveforms). ``$\Delta$ macro-$F_1$'' is the paired
  event-clustered bootstrap difference against SOPNet with its
  \SI{95}{\percent} percentile interval; all models are retrained on the
  identical training split with their official input protocols.}
  \label{tab:main}
  \small
  \begin{tabular}{lcccccc}
    \toprule
    Model & Accuracy & Macro-$F_1$ & Precision & Recall & MCC & $\Delta F_1$ [95\% CI] \\
    \midrule
    SOPNet (this work) & 0.9782 & 0.9780 & 0.9770 & 0.9754 & 0.9561 & --- \\
    Ross               & 0.9694 & 0.9692 & 0.9682 & 0.9650 & 0.9384 & $+0.0088$ [0.0083, 0.0094] \\
    DiTingMotion       & 0.9634 & 0.9631 & 0.9627 & 0.9573 & 0.9263 & $+0.0149$ [0.0143, 0.0155] \\
    RPNet              & 0.9728 & 0.9726 & 0.9708 & 0.9699 & 0.9453 & $+0.0054$ [0.0050, 0.0058] \\
    EQPolarity         & 0.9653 & 0.9651 & 0.9602 & 0.9643 & 0.9302 & $+0.0130$ [0.0124, 0.0135] \\
    CFM                & 0.9729 & 0.9727 & 0.9711 & 0.9697 & 0.9454 & $+0.0053$ [0.0049, 0.0057] \\
    \bottomrule
  \end{tabular}
\end{table}

\subsection{Per-source behaviour}
\label{sec:per_source}

Table~\ref{tab:source} breaks the strict accuracy down by source. SOPNet is
the most accurate model on every source (0.979 on SCSN, 0.980 on DiTing,
0.953 on PNW, 0.971 on INSTANCE and 0.587 on TXED). The three large sources
place the strongest models within about half a point of each other, whereas
the small sources carry visible sampling noise and spread the baselines
widely (INSTANCE 0.81--0.93, TXED 0.14--0.59 for known counts of 104 and 143
records, respectively). We report these numbers with their counts instead of
hiding them in an aggregate; the TXED behaviour is discussed in
Section~\ref{sec:discussion}.

\begin{table}[t]
  \centering
  \caption{Per-source strict accuracy on known-polarity test waveforms
  ($n = 253{,}194$ / $64{,}048$ / $9{,}742$ / $104$ / $143$ for SCSN /
  DiTing / PNW / INSTANCE / TXED, respectively).}
  \label{tab:source}
  \small
  \begin{tabular}{lccccc}
    \toprule
    Model & SCSN & DiTing & PNW & INSTANCE & TXED \\
    \midrule
    SOPNet       & 0.979 & 0.980 & 0.953 & 0.971 & 0.587 \\
    Ross         & 0.970 & 0.973 & 0.949 & 0.885 & 0.552 \\
    DiTingMotion & 0.965 & 0.973 & 0.870 & 0.808 & 0.140 \\
    RPNet        & 0.975 & 0.979 & 0.891 & 0.875 & 0.525 \\
    EQPolarity   & 0.969 & 0.973 & 0.837 & 0.827 & 0.168 \\
    CFM          & 0.974 & 0.978 & 0.913 & 0.933 & 0.399 \\
    \bottomrule
  \end{tabular}
\end{table}

\subsection{Ablation}
\label{sec:ablation}

To isolate where the signed field helps, we train four no-jitter variants under
identical data, schedule and loss scope: (A) SOPNet-Cls, a two-class
known-only classifier sharing the encoder; (B) the signed field with
known-only field loss and neither consistency term; (C) B plus the
polarity-consistency loss; (D) C plus the inversion-consistency loss (the full
model). Table~\ref{tab:ablation} reports their test metrics. The binary
classifier A (0.9746 accuracy / 0.9745 macro-$F_1$) and the unregularised
signed field B (0.9752 / 0.9751) are nearly indistinguishable, whereas the
full model D, which adds the polarity- and inversion-consistency losses,
improves to 0.9782 / 0.9780 ($+0.30$ points over B and $+0.36$ over A). The
advantage therefore comes from the consistency terms acting on the field
rather than from the field representation alone. Variant C (polarity loss
only) follows D's slow-convergence profile and was stopped at epoch~6 under a
scheduling constraint; its row is reported for completeness but is not
converged, so the loss attribution is based on the D-versus-B comparison.

\begin{table}[t]
  \centering
  \caption{Component ablation (no jitter; matched known-only loss scope; same
  backbone, data and schedule for all variants; test set, strict U/D).
  A: binary classification sharing the encoder; B: signed field; C: B plus
  the polarity-consistency loss; D: C plus the inversion-consistency loss
  (full model). C was stopped at epoch~6 under a scheduling constraint and is
  not converged (see text).}
  \label{tab:ablation}
  \small
  \begin{tabular}{lcccc}
    \toprule
    Variant & Accuracy & Macro-$F_1$ & MCC & Parameters \\
    \midrule
    A: SOPNet-Cls            & 0.9746 & 0.9745 & 0.9489 & 1.23M \\
    B: signed field          & 0.9752 & 0.9751 & 0.9501 & 1.44M \\
    C: $+$ polarity loss     & 0.9736 & 0.9734 & 0.9469 & 1.44M \\
    D: $+$ inversion (full)  & 0.9782 & 0.9780 & 0.9561 & 1.44M \\
    \bottomrule
  \end{tabular}
\end{table}

\subsection{Robustness}
\label{sec:robustness}

Figure~\ref{fig:snr} shows accuracy versus SNR for all six models on the fixed
subset. SOPNet is the most accurate model at every level and degrades
gracefully: 0.977 (clean), 0.963 (20\,dB), 0.938 (10\,dB), 0.909 (5\,dB),
0.862 (0\,dB) and 0.780 ($-5$\,dB); the baselines follow the same ordering
with gaps of 0.5--1.5 points. All models retain far-above-chance accuracy
even at the lowest SNR, consistent with first-motion information being
concentrated in the earliest samples of the onset.
Figure~\ref{fig:shift} shows the corresponding curves for P-reference shifts.
All six models degrade sharply beyond the zero reference: at $\pm$50\,ms
accuracies fall to 0.59--0.71, and at $\pm$100\,ms to 0.45--0.58, close to
chance. No model is robust; the collapse is a property of the fixed-window
classification regime rather than of any single architecture, and we return
to its implications in Section~\ref{sec:discussion}.

\begin{figure}[t]
  \centering
  \includegraphics[width=0.75\linewidth]{../../outputs/paper/figures/fig_snr_all.pdf}
  \caption{Strict up/down accuracy versus SNR for the six models on the fixed
  test subset. Curves are produced by \texttt{run\_robustness.py}.}
  \label{fig:snr}
\end{figure}

\begin{figure}[t]
  \centering
  \includegraphics[width=0.75\linewidth]{../../outputs/paper/figures/fig_p_shift_all.pdf}
  \caption{Strict up/down accuracy versus reference P-window shift.}
  \label{fig:shift}
\end{figure}

\begin{figure}[t]
  \centering
  \includegraphics[width=\linewidth]{../../outputs/paper/figures/fig5_calibration.pdf}
  \caption{Confidence calibration and selective prediction for SOPNet on
  known up/down test records: reliability diagram with and without Platt
  scaling (left) and accuracy as a function of retained coverage when
  thresholding the field-peak confidence (right).}
  \label{fig:calibration}
\end{figure}

\section{Discussion}
\label{sec:discussion}

\paragraph{Why a field?}
The ablation shows that the binary classifier (A, 0.9746 accuracy / 0.9745
macro-$F_1$) and the unregularised signed field (B, 0.9752 / 0.9751) are
nearly indistinguishable, whereas the full model (D, 0.9782 / 0.9780)
improves by 0.3--0.4 points over both. The benefit is therefore not the
field representation per se but what the consistency losses can do with it:
the attention-weighted polarity score and the inversion-equivariance
constraint shape the field's geometry, and only a full-resolution field ---
unlike a globally pooled classifier --- has the spatial resolution to carry
those constraints. The calibration results point the same way (D: ECE 0.031;
A: 0.038).

\paragraph{Known-only supervision.}
Excluding unknown-polarity records from the field loss was essential in
development: including their all-zero targets with full weight slows
convergence and blurs the onset lobe, while a down-weighted
(\texttt{unknown\_weight} $<1$) schedule leaves the final accuracy slightly
lower. This mirrors the label semantics: an unknown polarity is missing
information, not a third physical state, and treating it as a class forces the
model to model the annotator.

\paragraph{The P-shift trade-off.}
The fixed-centre window is the main limitation of the present implementation,
and Figure~\ref{fig:shift} makes the trade-off explicit: fixing the centre
lets the network devote its capacity to the onset waveform itself, yielding
high clean accuracy, but couples the decision to the accuracy of the
reference arrival. The six-model measurement sharpens the point: the collapse
at $\pm$50\,ms is regime-level, and the useful tolerance band is therefore
narrower than 50\,ms. Modern deep pickers on comparable regional data report
P-arrival errors of a few tens of milliseconds --- at the edge of this band.
Within the band the polarity readout is reliable; beyond it, two mitigations
apply. First, the failure mode is not silent: the field-peak magnitude drops
as the onset moves away from the trained centre, so the confidence threshold
(Section~\ref{sec:calibration}) can be used to abstain rather than to
mislabel. Second, part of the sensitivity is plausibly architectural: strided
downsampling makes encoder features sensitive to shifts that break the
sampling grid, which motivates shift-equivariant or anti-aliased
architectures. Training with random P jitter is the obvious data-side remedy
and also the obvious trap: jittering the field target forces the network to
solve localisation and polarity simultaneously, and slowed polarity
convergence by around seven accuracy points at short horizons in our
development runs. Removing the fixed centre altogether --- a learned
localisation objective --- is therefore the most promising direction for
future work, and we report the trade-off explicitly rather than optimising
one side of it.

\paragraph{Per-source variability.}
TXED's small known-polarity test subset (143 records) produces noisy and weak
estimates for every model, and the baselines span an unusually wide range
there (0.14--0.55 against SOPNet's 0.587); INSTANCE (104 records) separates
the models less widely but noisily (0.81--0.97). We report these numbers with
their counts instead of hiding them in an aggregate. The three large sources
show a consistent ordering, suggesting the benchmark is not dominated by a
single regional idiosyncrasy.

\paragraph{Limitations.}
The benchmark covers vertical-component, single-station records with
$\SI{100}{Hz}$ sampling and catalog P references; transfer to other
components, sampling rates or reference sources (for example, automated picks
with systematic bias) is not established here. The reported P-position error
is measured in the fixed-centre regime and should not be read as independent
picking performance. All results use a single training seed per model;
multi-seed repeats would strengthen the confidence in small differences and
are the first item of future validation work.

\section{Conclusion}
\label{sec:conclusion}

We presented SOPNet, a signed-onset-field network for P-wave first-motion
polarity determination, and a strict, event-level unified benchmark on five
open datasets with five retrained public baselines. Representing polarity as
the sign of a continuous field --- trained only on known labels and regularised
by polarity- and inversion-consistency losses --- matches the physics of the
first motion and yields state-of-the-art accuracy on the benchmark
(97.8\% on 327{,}231 known-polarity test waveforms) with calibrated
confidence and a tunable rejection option. Robustness experiments quantify
both the graceful degradation under noise and the sharp, previously
under-reported sensitivity to reference P shifts, which we frame as an
accuracy--tolerance trade-off of fixed-window classification. Extending the
field formulation to jointly refine the P arrival --- replacing the fixed
centre with a learned localisation objective --- is a natural next step toward
a fully self-contained polarity picker.

\section*{Data and code availability}
All five datasets are openly available (SCSN, DiTing, INSTANCE, PNW and
TXED). Code, configurations, trained checkpoints, the benchmark manifest and
all per-sample predictions will be released with the paper
\TBD{repository URL}.

\section*{Acknowledgments}
\TBD{}

\bibliographystyle{plainnat}
\bibliography{references}

\end{document}
